\section{Surface Model and Dipole Geometry}
\label{sec:surface_model}

We model the neutron star surface following the framework of
\citet{greenstein_hartke_1983}, in which anisotropic heat conduction through
the magnetised crust produces a latitude-dependent effective temperature.
The magnetic field is taken to be a centred dipole, tilted by an obliquity
angle $\alpha$ from the stellar spin axis.

\subsection{Magnetic colatitude}

At a geographic position $(\theta_{\mathrm{geo}},\,\varphi_{\mathrm{geo}})$
on the stellar surface, the magnetic colatitude $\theta_B$ with respect to
a dipole axis oriented along
$\hat{\boldsymbol{m}} = (\sin\alpha,\,0,\,\cos\alpha)$ is
\begin{equation}
  \cos\theta_B = \cos\theta_{\mathrm{geo}}\cos\alpha
    + \sin\theta_{\mathrm{geo}}\cos\varphi_{\mathrm{geo}}\sin\alpha\,.
  \label{eq:cos_thetaB}
\end{equation}
This is simply the dot product of the outward surface normal with the
magnetic axis unit vector, implemented as \texttt{magnetic\_colatitude()} in
\texttt{src/surface/dipole.jl}.

\subsection{Greenstein--Hartke temperature map}

\citet{greenstein_hartke_1983} show that the anisotropy of thermal
conductivity ($K_\parallel \gg K_\perp$) along versus across the magnetic
field lines channels heat preferentially toward the magnetic poles.  Their
Eq.~(14) gives the general relation
\begin{equation}
  T^{4} = T_{\mathrm{pole}}^{4}
    \!\left(\sin^{2}\Phi + \frac{K_\perp}{K_\parallel}\cos^{2}\Phi\right),
  \label{eq:GH_exact}
\end{equation}
where $\Phi$ is the angle between the magnetic field and the vertical,
related to the magnetic colatitude $\beta$ by
$\tan\Phi = 2\cot\beta$ (their Eq.~15).  In the limit
$K_\perp/K_\parallel \to 0$ the equatorial temperature vanishes; for
$K_\perp = K_\parallel$ the surface is isothermal.

Our implementation adopts the simplified parametric form used widely in the
literature:
\begin{equation}
  \boxed{T(\theta_B) = T_{\mathrm{eq}}
    + \bigl(T_{\mathrm{pole}} - T_{\mathrm{eq}}\bigr)\cos^{2}\!\theta_B\,,}
  \label{eq:T_surface}
\end{equation}
which recovers $T = T_{\mathrm{pole}}$ at the magnetic poles
($\theta_B = 0$) and $T = T_{\mathrm{eq}}$ at the magnetic equator
($\theta_B = \pi/2$).  This corresponds to the leading-order behaviour of
Eq.~\eqref{eq:GH_exact} when $T_{\mathrm{eq}}$ absorbs the residual
equatorial flux, and is cited as Eq.~(1) of \citet{greenstein_hartke_1983}
in the code documentation.  The function
\texttt{surface\_temperature($\theta_B$, $T_{\mathrm{pole}}$, $T_{\mathrm{eq}}$)}
enforces $T_{\mathrm{pole}} \geq T_{\mathrm{eq}} > 0$.

\subsection{Dipole magnetic field strength}

For a centred magnetic dipole, the surface field magnitude at magnetic
colatitude $\theta_B$ is the standard result
\begin{equation}
  \lvert\boldsymbol{B}\rvert(\theta_B)
    = \frac{B_{\mathrm{pole}}}{2}
      \sqrt{1 + 3\cos^{2}\!\theta_B}\,,
  \label{eq:B_dipole}
\end{equation}
giving $\lvert\boldsymbol{B}\rvert = B_{\mathrm{pole}}$ at the pole and
$\lvert\boldsymbol{B}\rvert = B_{\mathrm{pole}}/2$ at the equator
\citep{greenstein_hartke_1983}.  Both $T(\theta_B)$ and
$B(\theta_B)$ serve as inputs to the atmosphere grid lookup described in
Section~\ref{sec:atmosphere_grid}.


%======================================================================
\section{Schwarzschild Geodesic Ray Tracing}
\label{sec:ray_tracing}

Photons propagating outward from the neutron star surface travel along null
geodesics of the Schwarzschild metric.  We follow the analytic treatment of
\citet{beloborodov_2002} and cross-reference with the exact elliptic-integral
formulation of \citet{pechenick_1983}.

\subsection{The cosine relation}

Let a photon be emitted at the stellar surface $r = R$ at an angle $\alpha$
to the radial direction.  Observed at infinity, this photon escapes at a
polar angle $\psi$ from the sub-observer point.
\citet{beloborodov_2002}, Eq.~(1), gives the remarkably accurate approximate
relation
\begin{equation}
  \boxed{1 - \cos\alpha
    = (1 - \cos\psi)\!\left(1 - \frac{r_g}{R}\right),}
  \label{eq:beloborodov_cosine}
\end{equation}
where $r_g = 2GM/c^{2}$ is the Schwarzschild radius and
$u \equiv r_g/R$ is the compactness parameter.  This approximation is
valid for $R > 2r_g$ (i.e.\ $u < 1$) with fractional error $< 1\%$ for
the typical neutron star compactness $R \approx 3r_g$
\citep[Fig.~2]{beloborodov_2002}.

Inverting Eq.~\eqref{eq:beloborodov_cosine} yields
\begin{align}
  \cos\alpha &= 1 - (1 - \cos\psi)(1 - u)\,,
  \label{eq:cos_alpha} \\
  \cos\psi   &= 1 - \frac{1 - \cos\alpha}{1 - u}\,,
  \label{eq:cos_psi}
\end{align}
both implemented in \texttt{src/geodesics/schwarzschild.jl}.

\subsection{Impact parameter}

The impact parameter $b$ is related to the emission angle by
\citep[Eq.~10]{beloborodov_2002}
\begin{equation}
  \sin\alpha = \frac{b}{R}\sqrt{1 - \frac{r_g}{R}}\,,
  \label{eq:impact_param}
\end{equation}
so that
\begin{equation}
  b = \frac{R\sin\alpha}{\sqrt{1 - u}}\,.
  \label{eq:b_from_alpha}
\end{equation}
At the limb of the apparent disc the photon departs tangentially
($\cos\alpha = 0$, $\sin\alpha = 1$), giving the critical impact parameter
\begin{equation}
  b_{\mathrm{crit}} = \frac{R}{\sqrt{1 - u}}\,.
  \label{eq:b_crit}
\end{equation}
This agrees with Eq.~(2.10) of \citet{pechenick_1983},
$b_{\max} = R(1 - 2M/R)^{-1/2}$, noting the equivalence
$2M/R \equiv r_g/R = u$ in their geometrised units ($G = c = 1$).

\subsection{Exact bending-angle integral}

The exact bending angle is given by an elliptic integral over the photon
orbit \citep[Eq.~2.12]{pechenick_1983}:
\begin{equation}
  \Delta\phi
    = \int_{0}^{M/R}
      \left[u_b^{2} - (1 - 2u)\,u^{2}\right]^{-1/2}\!\mathrm{d}u\,,
  \label{eq:exact_bending}
\end{equation}
where $u_b \equiv M/b$ and the integration variable $u \equiv M/r$.  Our
code uses the Beloborodov approximation (Eq.~\ref{eq:beloborodov_cosine})
rather than evaluating this integral numerically, which introduces at most
$\sim 1\%$ error for $R \gtrsim 2r_g$ and avoids the need for special-function
libraries.

\subsection{Visible fraction of the stellar surface}

Gravitational light bending allows a distant observer to see more than the
Euclidean hemisphere.  The fraction of the surface visible is
\citep[Sect.~3]{beloborodov_2002}
\begin{equation}
  f_{\mathrm{vis}}
    = \frac{S_{v}}{4\pi R^{2}}
    = \frac{1}{2(1 - u)}\,.
  \label{eq:visible_fraction}
\end{equation}
In flat space ($u \to 0$), $f_{\mathrm{vis}} \to 1/2$ as expected.  For a
canonical neutron star with $R = 3\,r_g$ ($u = 2/3$),
$f_{\mathrm{vis}} = 3/2 \cdot 1/2 = 3/4$: three-quarters of the surface is
visible simultaneously.  The corresponding result in \citet{pechenick_1983}
follows from $(\Delta\phi)_{\max} = 180^{\circ}$ when $R/2M = 1.76$
(their p.~848), at which point the entire surface becomes visible.

\subsection{Image construction}

The \texttt{trace\_image()} routine constructs an $N \times N$ pixel grid in
the observer's image plane with Cartesian coordinates $(x,\,y)$ spanning
$[-b_{\max},\,b_{\max}]$ (where $b_{\max} = 1.05\,b_{\mathrm{crit}}$).
For each pixel:
\begin{enumerate}
  \item Compute the impact parameter $b = \sqrt{x^{2} + y^{2}}$ and image
    azimuth $\varphi_{\mathrm{img}} = \arctan(y/x)$.
  \item If $b > b_{\mathrm{crit}}$, the ray misses the star.
  \item Otherwise, obtain $\sin\alpha = (b/R)\sqrt{1 - u}$ from
    Eq.~\eqref{eq:impact_param}, then $\cos\alpha$ and $\cos\psi$ from
    Eqs.~\eqref{eq:cos_alpha}--\eqref{eq:cos_psi}.
  \item Map the sky position $(\psi,\,\varphi_{\mathrm{img}})$ to surface
    coordinates $(\theta_s,\,\varphi_s)$ via a spherical rotation accounting
    for the observer inclination $i$:
    \begin{align}
      x_s &= \sin i\cos\psi + \cos i\sin\psi\cos\varphi_{\mathrm{img}}\,,
        \notag \\
      y_s &= \sin\psi\sin\varphi_{\mathrm{img}}\,,
        \label{eq:sky_to_surface} \\
      z_s &= \cos i\cos\psi - \sin i\sin\psi\cos\varphi_{\mathrm{img}}\,,
        \notag
    \end{align}
    with $\theta_s = \arccos z_s$ and $\varphi_s = \arctan(y_s,\,x_s)$.
  \item Record the gravitational redshift $1 + z = (1 - u)^{-1/2}$,
    which is constant over the surface for the Schwarzschild metric.
\end{enumerate}
Each ray result stores
$\{b,\,\psi,\,\theta_s,\,\varphi_s,\,\cos\alpha,\,1{+}z\}$.


%======================================================================
\section{Relativistic Intensity Transform}
\label{sec:intensity_transform}

The specific intensity $I_\nu$ transforms under a gravitational redshift
according to the Lorentz-invariance of $I_\nu / \nu^{3}$
\citep[see e.g.][Sect.~3]{beloborodov_2002}:
\begin{equation}
  \frac{I_\nu^{\mathrm{obs}}}{\nu_{\mathrm{obs}}^{3}}
    = \frac{I_{\nu'}^{\mathrm{em}}}{\nu'^{3}}\,,
  \label{eq:invariant}
\end{equation}
where $\nu' = (1+z)\,\nu_{\mathrm{obs}}$ is the emitted frequency.
Rearranging,
\begin{equation}
  \boxed{I_\nu^{\mathrm{obs}}(\nu_{\mathrm{obs}})
    = \frac{I_{\nu'}^{\mathrm{em}}\!\bigl(\nu' = (1{+}z)\,\nu_{\mathrm{obs}}\bigr)}
           {(1 + z)^{3}}\,.}
  \label{eq:I_obs}
\end{equation}
This is the central equation applied per-pixel in the rendering pipeline.
The observed flux element from a surface patch $\mathrm{d}S$ at angle
$\alpha$ to the observer is \citep[Eq.~4]{beloborodov_2002}
\begin{equation}
  \mathrm{d}F = \left(1 - \frac{r_g}{R}\right)^{\!2}
    I_0(\alpha)\cos\alpha\,\frac{\mathrm{d}S}{D^{2}}\,,
  \label{eq:dF_beloborodov}
\end{equation}
where the $(1 - r_g/R)^{2}$ factor combines the $(1+z)^{-3}$ intensity
suppression with the $(1+z)^{+1}$ boost from the aberration of solid angles.
In our pixel-based rendering, each pixel subtends a known solid angle on
the sky, so we work directly with the specific intensity
(Eq.~\ref{eq:I_obs}) rather than integrating fluxes.


%======================================================================
\section{Spectral Image Cube and Atmosphere Grid Lookup}
\label{sec:atmosphere_grid}

\subsection{Pre-computed atmosphere grid}

To avoid re-solving the radiative transfer equation at every surface element,
we pre-compute atmosphere models on a two-dimensional grid in
$(T_{\mathrm{eff}},\,B)$ and store the emergent specific intensity
$I_\nu(\nu,\,\cos\theta_e)$ for each grid point
(\texttt{src/pipeline/atmosphere\_grid.jl}).  The \texttt{AtmosphereSpectrumGrid}
structure contains:
\begin{itemize}
  \item A common frequency grid $\{\nu_k\}_{k=1}^{K}$ and angle-cosine grid
    $\{\mu_j\}_{j=1}^{M}$;
  \item For each $(T_i,\,B_l)$: a $K \times M$ matrix of emergent intensities.
    For magnetic models ($B > 0$), the extraordinary (X) and ordinary (O)
    mode intensities are summed:
    $I_{\mathrm{total}} = I^{(X)} + I^{(O)}$.
\end{itemize}

\subsection{Grid interpolation}

At runtime, the lookup function
\texttt{lookup\_spectrum($T_{\mathrm{eff}}$, $B$, $\nu_{\mathrm{obs}}$,
$\cos\theta_e$)} performs trilinear interpolation:
\begin{enumerate}
  \item \emph{Bilinear in $(T_{\mathrm{eff}},\,B)$:} bracket the query
    values in the pre-computed grids and compute weights
    $w_{T}$, $w_{B}$ for the four surrounding models.
  \item \emph{Linear in angle:} for each of the four corner models,
    interpolate the $K \times M$ intensity table in the $\mu$-direction at
    $\mu = \cos\theta_e$.
  \item \emph{Nearest-neighbour in frequency:} for each requested
    $\nu_{\mathrm{obs}}$, the nearest grid frequency $\nu_k$ is used.
\end{enumerate}
The interpolated intensity at observed frequency $\nu_{\mathrm{obs}}$ is
thus
\begin{equation}
  I_{\nu}^{\mathrm{em}}(\nu',\,\cos\theta_e)
    \approx \sum_{a,b \in \{0,1\}} w_{T}^{(a)}\, w_{B}^{(b)}\,
    \Bigl[(1 - f_\mu)\, I_{k,\,j}^{(a,b)}
    + f_\mu\, I_{k,\,j+1}^{(a,b)}\Bigr],
  \label{eq:trilinear}
\end{equation}
where $f_\mu$ is the fractional position within the angle bracket and
$k$ is the nearest frequency index.

\subsection{Per-pixel spectrum computation}

The \texttt{render\_spectral\_cube()} function
(\texttt{src/pipeline/render.jl}) assembles the full spectral image cube
$I(x,\,y,\,\nu)$ via the following per-pixel procedure:
\begin{enumerate}
  \item \textbf{Ray trace} $\to$ surface coordinates
    $(\theta_s,\,\varphi_s)$, emission angle $\cos\alpha$, redshift $1{+}z$.
  \item \textbf{Surface model} $\to$ magnetic colatitude $\theta_B$
    (Eq.~\ref{eq:cos_thetaB}), local temperature $T(\theta_B)$
    (Eq.~\ref{eq:T_surface}), local field $B(\theta_B)$
    (Eq.~\ref{eq:B_dipole}).
  \item \textbf{Atmosphere lookup} $\to$ emitted spectrum
    $I_{\nu'}^{\mathrm{em}}(\nu_k,\,\cos\alpha)$ at surface-frame
    frequencies $\nu_k = (1{+}z)\,\nu_{\mathrm{obs},k}$
    via Eq.~\eqref{eq:trilinear}.
  \item \textbf{Redshift} $\to$ observed intensity per
    Eq.~\eqref{eq:I_obs}:
    \begin{equation}
      I_{\nu_k}^{\mathrm{obs}} = \frac{I_{\nu'}^{\mathrm{em}}(\nu_k)}{(1 + z)^{3}}\,.
      \label{eq:pixel_redshift}
    \end{equation}
\end{enumerate}
The result is a data cube
$\texttt{SpectralImageCube}(N_x,\,N_y,\,N_\nu)$ storing $I(x,y,\nu)$
in units of $\mathrm{erg\,s^{-1}\,cm^{-2}\,Hz^{-1}\,sr^{-1}}$, together
with per-pixel metadata (hit fraction, local $T$, $B$, and $1{+}z$).


%======================================================================
\section{CIE~1931 Colorimetry and sRGB Rendering}
\label{sec:colorimetry}

The spectral image cube must be converted to a displayable colour image.
We implement the standard colorimetric pipeline in
\texttt{src/colorimetry/cie\_srgb.jl}.

\subsection{CIE~1931 XYZ tristimulus values}

The CIE~1931 $2^{\circ}$ colour matching functions (CMFs)
$\bar{x}(\lambda)$, $\bar{y}(\lambda)$, $\bar{z}(\lambda)$ are loaded
from the CVRL dataset (360--830\,nm at 1\,nm resolution).  The tristimulus
values for an observed spectrum $I_\nu(\nu)$ are obtained by transforming to
wavelength space and integrating against the CMFs:
\begin{equation}
  X = \int I_\lambda(\lambda)\,\bar{x}(\lambda)\,\mathrm{d}\lambda\,,
  \qquad
  Y = \int I_\lambda(\lambda)\,\bar{y}(\lambda)\,\mathrm{d}\lambda\,,
  \qquad
  Z = \int I_\lambda(\lambda)\,\bar{z}(\lambda)\,\mathrm{d}\lambda\,,
  \label{eq:XYZ}
\end{equation}
where the change of variables from frequency to wavelength gives
\begin{equation}
  I_\lambda = I_\nu \,\frac{c}{\lambda^{2}}\,,
  \label{eq:I_lambda}
\end{equation}
using $|\mathrm{d}\nu| = (c/\lambda^{2})\,|\mathrm{d}\lambda|$.
The integrals are evaluated by trapezoidal quadrature over the CMF
wavelength grid, with $I_\nu$ interpolated in log-log space (binary search
plus linear interpolation of $\ln I_\nu$ vs.\ $\ln\nu$) to handle the
large dynamic range of astrophysical spectra.

\subsection{XYZ to linear sRGB}

The CIE~XYZ tristimulus values are converted to linear sRGB using the
IEC~61966-2-1:1999 matrix for the D65 illuminant:
\begin{equation}
  \begin{pmatrix} R_{\mathrm{lin}} \\ G_{\mathrm{lin}} \\ B_{\mathrm{lin}} \end{pmatrix}
  = \begin{pmatrix}
    \phantom{-}3.2406 & -1.5372 & -0.4986 \\
    -0.9689 & \phantom{-}1.8758 & \phantom{-}0.0415 \\
    \phantom{-}0.0557 & -0.2040 & \phantom{-}1.0570
  \end{pmatrix}
  \begin{pmatrix} X \\ Y \\ Z \end{pmatrix}.
  \label{eq:XYZ_to_sRGB}
\end{equation}
Negative values (out-of-gamut colours) are clamped to zero.

\subsection{Tone mapping}

Neutron star spectra produce luminance values spanning many orders of
magnitude.  We apply the global Reinhard tone mapping operator
\citep{reinhard_2002}, which compresses the luminance channel $Y$ via
\begin{equation}
  L_{\mathrm{scaled}} = \frac{a \cdot Y}{\bar{Y}}\,,
  \qquad
  Y_{\mathrm{display}} = \frac{L_{\mathrm{scaled}}}{1 + L_{\mathrm{scaled}}}\,,
  \label{eq:reinhard}
\end{equation}
where $\bar{Y}$ is the geometric mean luminance of all illuminated pixels
and $a = 0.18$ is the key value (scene middle grey).  The chromaticity is
preserved by scaling all three XYZ channels by
$Y_{\mathrm{display}} / Y$ before the matrix conversion
(Eq.~\ref{eq:XYZ_to_sRGB}).

\subsection{sRGB gamma correction}

The final step applies the piecewise sRGB transfer function
(IEC~61966-2-1):
\begin{equation}
  C_{\mathrm{sRGB}} =
  \begin{cases}
    12.92\, C_{\mathrm{lin}}\,, & C_{\mathrm{lin}} \leq 0.0031308\,, \\[4pt]
    1.055\, C_{\mathrm{lin}}^{1/2.4} - 0.055\,, & C_{\mathrm{lin}} > 0.0031308\,,
  \end{cases}
  \label{eq:srgb_gamma}
\end{equation}
applied independently to each of $R$, $G$, $B$.  The result is quantised
to 8-bit integers $\in [0,\,255]$ and written as a PPM image.

\subsection{False-colour X-ray rendering}

In addition to the true-colour rendering, the pipeline produces a
false-colour X-ray image by integrating the spectral cube over three
energy bands and mapping them to RGB channels:
\begin{equation}
  \text{Red}: 0.1\text{--}0.5~\mathrm{keV}\,,\quad
  \text{Green}: 0.5\text{--}2.0~\mathrm{keV}\,,\quad
  \text{Blue}: 2.0\text{--}10.0~\mathrm{keV}\,.
  \label{eq:false_colour_bands}
\end{equation}
Each channel is independently normalised to $[0,\,255]$, providing a visual
diagnostic of spectral hardness variations across the stellar surface.